\documentclass[11pt,a4paper]{article}

%=========================================================
% PACKAGES
%=========================================================
\usepackage[utf8]{inputenc}
\usepackage[T1]{fontenc}
\usepackage[french]{babel}
\usepackage[a4paper,margin=2.1cm]{geometry}
\usepackage{lmodern}
\usepackage{amsmath,amssymb,amsthm,mathtools}
\usepackage{fancyhdr}
\usepackage{lastpage}
\usepackage{enumitem}
\usepackage{array,tabularx}
\usepackage{tikz}
\usepackage[most]{tcolorbox}
\usepackage{hyperref}
\usepackage{xcolor}

%=========================================================
% MISE EN PAGE
%=========================================================
\setlength{\headheight}{14pt}
\pagestyle{fancy}
\fancyhf{}
\fancyhead[L]{Analyse CPGE}
\fancyhead[C]{Fonctions réelles}
\fancyhead[R]{Page \thepage/\pageref{LastPage}}
\fancyfoot[C]{Limites -- Continuité -- Bijections}

\hypersetup{
    colorlinks=true,
    linkcolor=blue!50!black,
    urlcolor=blue!50!black
}

\setlist[itemize]{label=\(\triangleright\)}
\setlist[enumerate]{label=\textbf{\arabic*.}}

%=========================================================
% COULEURS ET BOÎTES
%=========================================================
\definecolor{bleufonce}{RGB}{25,65,130}
\definecolor{bleuclair}{RGB}{235,243,255}
\definecolor{vertfonce}{RGB}{30,120,70}
\definecolor{vertclair}{RGB}{235,250,240}
\definecolor{rougefonce}{RGB}{170,40,40}
\definecolor{rougeclair}{RGB}{255,238,238}
\definecolor{orangefonce}{RGB}{190,110,20}
\definecolor{orangeclair}{RGB}{255,246,230}
\definecolor{violetfonce}{RGB}{95,55,145}
\definecolor{violetclair}{RGB}{246,241,255}

\tcbset{
    enhanced,
    sharp corners,
    boxrule=0.75pt,
    before skip=8pt,
    after skip=8pt
}

\newtcolorbox{definitionbox}[1][]{
    colback=bleuclair,
    colframe=bleufonce,
    title=\textbf{Définition #1}
}

\newtcolorbox{theorembox}[1][]{
    colback=vertclair,
    colframe=vertfonce,
    title=\textbf{Théorème / Propriété #1}
}

\newtcolorbox{methodbox}[1][]{
    colback=orangeclair,
    colframe=orangefonce,
    title=\textbf{Méthode #1}
}

\newtcolorbox{retenirbox}[1][]{
    colback=violetclair,
    colframe=violetfonce,
    title=\textbf{À retenir #1}
}

\newtcolorbox{erreurbox}[1][]{
    colback=rougeclair,
    colframe=rougefonce,
    title=\textbf{Erreur classique #1}
}

\newtcolorbox{exercicebox}[1][]{
    colback=white,
    colframe=black!60,
    title=\textbf{Exercice #1}
}

\newtcolorbox{correctionbox}[1][]{
    colback=gray!5,
    colframe=black!65,
    title=\textbf{Correction #1}
}

%=========================================================
% COMMANDES
%=========================================================
\newcommand{\N}{\mathbb N}
\newcommand{\Z}{\mathbb Z}
\newcommand{\Q}{\mathbb Q}
\newcommand{\R}{\mathbb R}
\newcommand{\C}{\mathbb C}
\newcommand{\abs}[1]{\left|#1\right|}
\newcommand{\ssi}{\Longleftrightarrow}
\newcommand{\implique}{\Longrightarrow}
\newcommand{\limx}[2]{\lim_{#1\to #2}}
\newcommand{\Df}{\mathcal D_f}
\newcommand{\Id}{\operatorname{Id}}
\newcommand{\Imf}{\operatorname{Im}}
\newcommand{\Card}{\operatorname{Card}}

\newtheorem{prop}{Proposition}[section]
\newtheorem{lemme}[prop]{Lemme}
\newtheorem{corollaire}[prop]{Corollaire}

%=========================================================
% DOCUMENT
%=========================================================
\begin{document}

\begin{center}
    {\Huge \textbf{Fonctions réelles}}\\[2mm]
    {\Large Limites, continuité et bijections}\\[2mm]
    {\large Cours complet pour classes préparatoires scientifiques}\\[2mm]
    \rule{0.88\linewidth}{0.5pt}
\end{center}

\vspace{0.4cm}

\tableofcontents

\newpage

%=========================================================
\section*{Introduction générale}
\addcontentsline{toc}{section}{Introduction générale}

Après les suites numériques, on étudie les fonctions réelles d'une variable réelle.  
Une fonction réelle permet de modéliser une quantité dépendant continûment ou non d'une variable réelle.

Les notions centrales de ce chapitre sont :
\begin{itemize}
    \item le domaine de définition ;
    \item l'image directe et l'image réciproque ;
    \item la limite d'une fonction ;
    \item les limites à gauche et à droite ;
    \item la continuité ;
    \item le prolongement par continuité ;
    \item le théorème des valeurs intermédiaires ;
    \item le théorème de la bijection continue ;
    \item la fonction réciproque.
\end{itemize}

\begin{retenirbox}
Les suites étudient le comportement d'une quantité lorsque l'indice entier \(n\) tend vers \(+\infty\).  
Les fonctions réelles permettent d'étudier le comportement d'une quantité lorsque la variable réelle \(x\) tend vers un réel \(a\), vers \(+\infty\), ou vers \(-\infty\).
\end{retenirbox}

Ce chapitre est fondamental pour la suite du cours :
\begin{itemize}
    \item la dérivation repose sur la limite d'un taux d'accroissement ;
    \item l'intégration utilise la continuité ;
    \item les équations différentielles utilisent les fonctions continues ;
    \item les développements limités sont des approximations locales de fonctions ;
    \item le calcul différentiel généralise ces notions à plusieurs variables.
\end{itemize}

%=========================================================
\section{Fonctions réelles d'une variable réelle}

\subsection{Définition générale}

\begin{definitionbox}[Fonction réelle d'une variable réelle]
Une fonction réelle d'une variable réelle est une application :
\[
f:D\longrightarrow \R,
\]
où \(D\) est une partie de \(\R\).  
À chaque réel \(x\in D\), elle associe un unique réel \(f(x)\).
\end{definitionbox}

On écrit :
\[
x\longmapsto f(x).
\]

Le réel \(x\) est appelé la variable.  
Le réel \(f(x)\) est appelé l'image de \(x\) par \(f\).

\begin{retenirbox}
Une fonction est définie par trois données :
\begin{itemize}
    \item un ensemble de départ ;
    \item un ensemble d'arrivée ;
    \item une règle d'association.
\end{itemize}
En analyse réelle, on travaille souvent avec :
\[
f:D\subset\R\longrightarrow \R.
\]
\end{retenirbox}

\subsection{Domaine de définition}

\begin{definitionbox}[Domaine de définition]
Le domaine de définition d'une fonction \(f\), noté \(\Df\), est l'ensemble des réels \(x\) pour lesquels l'expression \(f(x)\) a un sens.
\end{definitionbox}

\begin{methodbox}[Déterminer un domaine de définition]
Pour déterminer le domaine de définition d'une fonction réelle, on vérifie :
\begin{enumerate}
    \item les dénominateurs ne doivent pas être nuls ;
    \item les quantités sous une racine carrée doivent être positives ou nulles ;
    \item les arguments de logarithmes doivent être strictement positifs ;
    \item les fonctions trigonométriques réciproques ont des domaines restreints ;
    \item les compositions doivent être cohérentes.
\end{enumerate}
\end{methodbox}

\begin{exercicebox}[Domaine de définition]
Déterminer le domaine de définition de :
\[
f(x)=\frac{\sqrt{x+1}}{x-2}.
\]
\end{exercicebox}

\begin{correctionbox}
On doit avoir :
\[
x+1\geq 0,
\]
c'est-à-dire :
\[
x\geq -1.
\]
On doit aussi avoir :
\[
x-2\neq0,
\]
c'est-à-dire :
\[
x\neq2.
\]
Donc :
\[
\Df=[-1,+\infty[\setminus\{2\}.
\]
\end{correctionbox}

\subsection{Graphe d'une fonction}

\begin{definitionbox}[Graphe]
Le graphe de \(f:D\to\R\) est l'ensemble :
\[
\mathcal G_f=\{(x,f(x)):\ x\in D\}\subset\R^2.
\]
\end{definitionbox}

Le graphe est un objet géométrique, mais les propriétés d'une fonction doivent être démontrées par des arguments mathématiques, et non seulement par lecture graphique.

\begin{erreurbox}
Un dessin permet de conjecturer, mais ne constitue pas une preuve.  
Par exemple, une fonction peut sembler croissante sur un graphique sans l'être réellement.
\end{erreurbox}

%=========================================================
\section{Image directe et image réciproque}

\subsection{Image directe}

\begin{definitionbox}[Image directe]
Soit \(f:D\to\R\), et soit \(A\subset D\).  
L'image directe de \(A\) par \(f\) est :
\[
f(A)=\{f(x):x\in A\}.
\]
\end{definitionbox}

Autrement dit :
\[
y\in f(A)
\ssi
\exists x\in A,\quad y=f(x).
\]

\begin{exercicebox}[Image directe]
Soit :
\[
f(x)=x^2.
\]
Déterminer :
\[
f([-2,3]).
\]
\end{exercicebox}

\begin{correctionbox}
Sur l'intervalle \([-2,3]\), la fonction \(x\mapsto x^2\) atteint son minimum en \(0\), qui appartient à l'intervalle.  
On a :
\[
f(0)=0.
\]
Le maximum sur \([-2,3]\) est atteint en \(3\), car :
\[
3^2=9
\]
et :
\[
(-2)^2=4.
\]
Donc :
\[
f([-2,3])=[0,9].
\]
\end{correctionbox}

\subsection{Image réciproque}

\begin{definitionbox}[Image réciproque]
Soit \(f:D\to\R\), et soit \(B\subset\R\).  
L'image réciproque de \(B\) par \(f\) est :
\[
f^{-1}(B)=\{x\in D:\ f(x)\in B\}.
\]
\end{definitionbox}

Attention : cette notation ne suppose pas que \(f\) soit bijective.

\begin{retenirbox}
Pour toute fonction \(f\), l'écriture :
\[
f^{-1}(B)
\]
désigne l'image réciproque d'un ensemble \(B\).  
Elle existe même si \(f\) n'a pas de fonction réciproque.
\end{retenirbox}

\begin{exercicebox}[Image réciproque]
Soit :
\[
f(x)=x^2.
\]
Déterminer :
\[
f^{-1}([1,4]).
\]
\end{exercicebox}

\begin{correctionbox}
On cherche les réels \(x\) tels que :
\[
x^2\in[1,4].
\]
Cela équivaut à :
\[
1\leq x^2\leq 4.
\]
Donc :
\[
1\leq |x|\leq2.
\]
Ainsi :
\[
x\in[-2,-1]\cup[1,2].
\]
Donc :
\[
f^{-1}([1,4])=[-2,-1]\cup[1,2].
\]
\end{correctionbox}

%=========================================================
\section{Parité et périodicité}

\subsection{Fonctions paires et impaires}

\begin{definitionbox}[Ensemble symétrique par rapport à \(0\)]
Une partie \(D\subset\R\) est symétrique par rapport à \(0\) si :
\[
\forall x\in D,\quad -x\in D.
\]
\end{definitionbox}

\begin{definitionbox}[Fonction paire]
Soit \(f:D\to\R\), où \(D\) est symétrique par rapport à \(0\).  
La fonction \(f\) est paire si :
\[
\forall x\in D,\quad f(-x)=f(x).
\]
\end{definitionbox}

\begin{definitionbox}[Fonction impaire]
Soit \(f:D\to\R\), où \(D\) est symétrique par rapport à \(0\).  
La fonction \(f\) est impaire si :
\[
\forall x\in D,\quad f(-x)=-f(x).
\]
\end{definitionbox}

\begin{retenirbox}
Une fonction paire a un graphe symétrique par rapport à l'axe des ordonnées.  
Une fonction impaire a un graphe symétrique par rapport à l'origine.
\end{retenirbox}

\begin{erreurbox}
Avant de parler de parité, il faut vérifier que le domaine de définition est symétrique par rapport à \(0\).
\end{erreurbox}

\begin{exercicebox}[Parité]
Étudier la parité de :
\[
f(x)=x^4-3x^2+1.
\]
\end{exercicebox}

\begin{correctionbox}
Le domaine de définition est \(\R\), qui est symétrique par rapport à \(0\).  
On calcule :
\[
f(-x)=(-x)^4-3(-x)^2+1=x^4-3x^2+1=f(x).
\]
Donc \(f\) est paire.
\end{correctionbox}

\subsection{Périodicité}

\begin{definitionbox}[Fonction périodique]
Soit \(f:D\to\R\).  
La fonction \(f\) est périodique s'il existe \(T>0\) tel que :
\[
\forall x\in D,\quad x+T\in D
\]
et :
\[
f(x+T)=f(x).
\]
Un tel réel \(T\) est appelé période de \(f\).
\end{definitionbox}

\begin{retenirbox}
Si \(T\) est une période de \(f\), alors tout multiple entier non nul de \(T\) est aussi une période.
\end{retenirbox}

\begin{exercicebox}[Périodicité]
Montrer que la fonction :
\[
f(x)=\sin(3x)
\]
est périodique et déterminer une période.
\end{exercicebox}

\begin{correctionbox}
On sait que :
\[
\sin(t+2\pi)=\sin(t).
\]
On cherche \(T>0\) tel que :
\[
3(x+T)=3x+2\pi.
\]
Cela donne :
\[
3T=2\pi,
\]
donc :
\[
T=\frac{2\pi}{3}.
\]
Alors :
\[
f\left(x+\frac{2\pi}{3}\right)
=
\sin\left(3x+2\pi\right)
=
\sin(3x)
=
f(x).
\]
Donc \(f\) est périodique de période \(\frac{2\pi}{3}\).
\end{correctionbox}

%=========================================================
\section{Limite finie d'une fonction en un point}

\subsection{Définition}

\begin{definitionbox}[Limite finie en un point]
Soit \(f:D\to\R\), soit \(a\in\R\), et soit \(\ell\in\R\).  
On dit que \(f(x)\) tend vers \(\ell\) lorsque \(x\) tend vers \(a\), avec \(x\in D\), si :
\[
\forall \varepsilon>0,\quad \exists \eta>0,\quad \forall x\in D,\quad
0<|x-a|\leq \eta \implique |f(x)-\ell|\leq\varepsilon.
\]
On note :
\[
\lim_{x\to a}f(x)=\ell.
\]
\end{definitionbox}

\begin{retenirbox}
La condition :
\[
0<|x-a|
\]
signifie que l'on regarde les valeurs de \(f(x)\) lorsque \(x\) est proche de \(a\), mais éventuellement différent de \(a\).  
La valeur de \(f(a)\), si elle existe, n'intervient pas dans la limite.
\end{retenirbox}

\subsection{Interprétation}

La définition :
\[
\forall \varepsilon>0,\quad \exists \eta>0,\quad
0<|x-a|\leq\eta\implique |f(x)-\ell|\leq\varepsilon
\]
signifie que l'on peut rendre \(f(x)\) aussi proche que l'on veut de \(\ell\), à condition de prendre \(x\) suffisamment proche de \(a\), mais distinct de \(a\).

\begin{erreurbox}
Ne pas confondre :
\[
\lim_{x\to a}f(x)
\]
et :
\[
f(a).
\]
Une fonction peut avoir une limite en \(a\) sans être définie en \(a\).
\end{erreurbox}

\subsection{Exemple avec la définition}

\begin{exercicebox}[Limite avec la définition]
Montrer que :
\[
\lim_{x\to 2}(3x+1)=7.
\]
\end{exercicebox}

\begin{correctionbox}
On veut montrer :
\[
\forall\varepsilon>0,\quad \exists\eta>0,\quad |x-2|\leq\eta\implique |(3x+1)-7|\leq\varepsilon.
\]
On calcule :
\[
|(3x+1)-7|=|3x-6|=3|x-2|.
\]
Il suffit donc d'imposer :
\[
3|x-2|\leq\varepsilon,
\]
c'est-à-dire :
\[
|x-2|\leq\frac{\varepsilon}{3}.
\]
On choisit :
\[
\eta=\frac{\varepsilon}{3}.
\]
Alors :
\[
|x-2|\leq\eta
\implique
|(3x+1)-7|\leq\varepsilon.
\]
Donc :
\[
\lim_{x\to2}(3x+1)=7.
\]
\end{correctionbox}

%=========================================================
\section{Limites infinies et limites à l'infini}

\subsection{Limite infinie en un point}

\begin{definitionbox}[Limite \(+\infty\) en un point]
Soit \(f:D\to\R\) et \(a\in\R\).  
On dit que :
\[
\lim_{x\to a}f(x)=+\infty
\]
si :
\[
\forall A\in\R,\quad \exists\eta>0,\quad \forall x\in D,\quad
0<|x-a|\leq\eta\implique f(x)\geq A.
\]
\end{definitionbox}

\begin{definitionbox}[Limite \(-\infty\) en un point]
On dit que :
\[
\lim_{x\to a}f(x)=-\infty
\]
si :
\[
\forall A\in\R,\quad \exists\eta>0,\quad \forall x\in D,\quad
0<|x-a|\leq\eta\implique f(x)\leq A.
\]
\end{definitionbox}

\subsection{Limite finie à l'infini}

\begin{definitionbox}[Limite finie en \(+\infty\)]
Soit \(f:D\to\R\) et \(\ell\in\R\).  
On dit que :
\[
\lim_{x\to+\infty}f(x)=\ell
\]
si :
\[
\forall \varepsilon>0,\quad \exists A\in\R,\quad \forall x\in D,\quad
x\geq A\implique |f(x)-\ell|\leq\varepsilon.
\]
\end{definitionbox}

\begin{definitionbox}[Limite \(+\infty\) en \(+\infty\)]
On dit que :
\[
\lim_{x\to+\infty}f(x)=+\infty
\]
si :
\[
\forall B\in\R,\quad \exists A\in\R,\quad \forall x\in D,\quad
x\geq A\implique f(x)\geq B.
\]
\end{definitionbox}

Les définitions en \(-\infty\) sont analogues.

\begin{retenirbox}
On rencontre quatre grands types de limites :
\[
x\to a,\quad x\to+\infty,\quad x\to-\infty,
\]
et la limite peut être :
\[
\ell\in\R,\quad +\infty,\quad -\infty.
\]
\end{retenirbox}

%=========================================================
\section{Limites à gauche et à droite}

\subsection{Définition}

\begin{definitionbox}[Limite à droite]
Soit \(f:D\to\R\), \(a\in\R\), et \(\ell\in\R\).  
On dit que \(f\) admet \(\ell\) pour limite à droite en \(a\) si :
\[
\forall \varepsilon>0,\quad \exists \eta>0,\quad \forall x\in D,\quad
a<x\leq a+\eta \implique |f(x)-\ell|\leq\varepsilon.
\]
On note :
\[
\lim_{x\to a^+}f(x)=\ell.
\]
\end{definitionbox}

\begin{definitionbox}[Limite à gauche]
On dit que \(f\) admet \(\ell\) pour limite à gauche en \(a\) si :
\[
\forall \varepsilon>0,\quad \exists \eta>0,\quad \forall x\in D,\quad
a-\eta\leq x<a \implique |f(x)-\ell|\leq\varepsilon.
\]
On note :
\[
\lim_{x\to a^-}f(x)=\ell.
\]
\end{definitionbox}

\subsection{Lien avec la limite bilatérale}

\begin{theorembox}[Limite bilatérale et limites latérales]
Soit \(f\) définie au voisinage de \(a\), éventuellement privée de \(a\).  
Alors :
\[
\lim_{x\to a}f(x)=\ell
\]
si et seulement si :
\[
\lim_{x\to a^-}f(x)=\ell
\quad\text{et}\quad
\lim_{x\to a^+}f(x)=\ell.
\]
\end{theorembox}

\begin{proof}
Si \(f(x)\to\ell\) lorsque \(x\to a\), alors la condition est vraie en particulier lorsque \(x<a\) et lorsque \(x>a\).  
Les limites à gauche et à droite valent donc \(\ell\).

Réciproquement, supposons :
\[
\lim_{x\to a^-}f(x)=\ell
\quad\text{et}\quad
\lim_{x\to a^+}f(x)=\ell.
\]
Soit \(\varepsilon>0\).  
Il existe \(\eta_1>0\) tel que :
\[
a-\eta_1\leq x<a\implique |f(x)-\ell|\leq\varepsilon.
\]
Il existe \(\eta_2>0\) tel que :
\[
a<x\leq a+\eta_2\implique |f(x)-\ell|\leq\varepsilon.
\]
Posons :
\[
\eta=\min(\eta_1,\eta_2).
\]
Alors, si :
\[
0<|x-a|\leq\eta,
\]
on a soit \(x<a\), soit \(x>a\), et dans les deux cas :
\[
|f(x)-\ell|\leq\varepsilon.
\]
Donc :
\[
\lim_{x\to a}f(x)=\ell.
\]
\end{proof}

\begin{erreurbox}
Si les limites à gauche et à droite existent mais sont différentes, alors la limite en \(a\) n'existe pas.
\end{erreurbox}

\begin{exercicebox}[Limites latérales]
Soit :
\[
f(x)=
\begin{cases}
0 & \text{si } x<0,\\
1 & \text{si } x>0.
\end{cases}
\]
Étudier la limite de \(f\) en \(0\).
\end{exercicebox}

\begin{correctionbox}
On a :
\[
\lim_{x\to0^-}f(x)=0,
\]
car pour \(x<0\), \(f(x)=0\).  
Et :
\[
\lim_{x\to0^+}f(x)=1,
\]
car pour \(x>0\), \(f(x)=1\).  
Les deux limites latérales sont différentes.  
Donc :
\[
\lim_{x\to0}f(x)
\]
n'existe pas.
\end{correctionbox}

%=========================================================
\section{Opérations sur les limites}

\subsection{Somme, produit, quotient}

\begin{theorembox}[Opérations sur les limites]
Soient \(f\) et \(g\) deux fonctions définies au voisinage de \(a\), éventuellement privé de \(a\).  
Si :
\[
\lim_{x\to a}f(x)=\ell
\quad\text{et}\quad
\lim_{x\to a}g(x)=m,
\]
avec \(\ell,m\in\R\), alors :
\[
\lim_{x\to a}(f(x)+g(x))=\ell+m,
\]
\[
\lim_{x\to a}(f(x)g(x))=\ell m.
\]
Si \(m\neq0\) et \(g(x)\neq0\) au voisinage de \(a\), alors :
\[
\lim_{x\to a}\frac{f(x)}{g(x)}=\frac{\ell}{m}.
\]
\end{theorembox}

\begin{proof}
Les démonstrations sont analogues à celles des suites.  
Pour la somme, on écrit :
\[
|(f(x)+g(x))-(\ell+m)|
\leq |f(x)-\ell|+|g(x)-m|.
\]
Pour le produit :
\[
f(x)g(x)-\ell m=f(x)(g(x)-m)+m(f(x)-\ell),
\]
et on utilise le fait qu'une fonction admettant une limite finie est bornée au voisinage du point considéré.  
Pour le quotient, on montre d'abord que :
\[
\frac1{g(x)}\to\frac1m
\]
en utilisant que \(g(x)\) reste éloignée de \(0\) au voisinage de \(a\).
\end{proof}

\subsection{Composition}

\begin{theorembox}[Limite d'une composée]
Soient \(f\) et \(g\) deux fonctions.  
Si :
\[
\lim_{x\to a}g(x)=b
\]
et :
\[
\lim_{y\to b}f(y)=\ell,
\]
alors, sous les conditions usuelles de définition :
\[
\lim_{x\to a}f(g(x))=\ell.
\]
\end{theorembox}

\begin{proof}
Soit \(\varepsilon>0\).  
Comme :
\[
\lim_{y\to b}f(y)=\ell,
\]
il existe \(\alpha>0\) tel que :
\[
0<|y-b|\leq\alpha \implique |f(y)-\ell|\leq\varepsilon.
\]
Comme :
\[
g(x)\to b,
\]
il existe \(\eta>0\) tel que :
\[
0<|x-a|\leq\eta \implique |g(x)-b|\leq\alpha.
\]
Alors, pour \(x\) suffisamment proche de \(a\), \(g(x)\) est suffisamment proche de \(b\), donc \(f(g(x))\) est proche de \(\ell\).  
On obtient :
\[
f(g(x))\to\ell.
\]
\end{proof}

\begin{erreurbox}
La composition impose de vérifier que les expressions sont bien définies.  
Il ne suffit pas d'écrire mécaniquement :
\[
\lim f(g(x))=f(\lim g(x)).
\]
Cette écriture n'est valable que si \(f\) est continue en la limite de \(g\).
\end{erreurbox}

%=========================================================
\section{Limites et suites : caractérisation séquentielle}

\begin{theorembox}[Caractérisation séquentielle de la limite]
Soit \(f:D\to\R\), soit \(a\in\R\), et soit \(\ell\in\R\).  
Alors :
\[
\lim_{x\to a}f(x)=\ell
\]
si et seulement si, pour toute suite \((x_n)\) d'éléments de \(D\setminus\{a\}\) telle que :
\[
x_n\to a,
\]
on a :
\[
f(x_n)\to\ell.
\]
\end{theorembox}

\begin{proof}
Supposons :
\[
\lim_{x\to a}f(x)=\ell.
\]
Soit \((x_n)\) une suite d'éléments de \(D\setminus\{a\}\) telle que :
\[
x_n\to a.
\]
Soit \(\varepsilon>0\).  
Par définition de la limite, il existe \(\eta>0\) tel que :
\[
0<|x-a|\leq\eta \implique |f(x)-\ell|\leq\varepsilon.
\]
Comme :
\[
x_n\to a,
\]
il existe \(N\in\N\) tel que :
\[
n\geq N\implique |x_n-a|\leq\eta.
\]
Comme \(x_n\neq a\), on a :
\[
0<|x_n-a|\leq\eta.
\]
Donc :
\[
|f(x_n)-\ell|\leq\varepsilon.
\]
Ainsi :
\[
f(x_n)\to\ell.
\]

Réciproquement, on démontre la contraposée.  
Supposons que \(f(x)\) ne tende pas vers \(\ell\) lorsque \(x\to a\).  
Alors il existe \(\varepsilon_0>0\) tel que, pour tout \(\eta>0\), il existe \(x\in D\) vérifiant :
\[
0<|x-a|\leq\eta
\]
et :
\[
|f(x)-\ell|>\varepsilon_0.
\]
Pour chaque \(n\geq1\), on prend \(\eta=\frac1n\).  
Il existe donc \(x_n\in D\) tel que :
\[
0<|x_n-a|\leq\frac1n
\]
et :
\[
|f(x_n)-\ell|>\varepsilon_0.
\]
Alors :
\[
x_n\to a,
\]
mais \(f(x_n)\) ne tend pas vers \(\ell\).  
Cela contredit la propriété séquentielle.
\end{proof}

\begin{methodbox}[Prouver qu'une limite n'existe pas]
Pour montrer que \(f\) n'admet pas de limite en \(a\), on peut trouver deux suites :
\[
x_n\to a,\qquad y_n\to a,
\]
avec \(x_n\neq a\), \(y_n\neq a\), telles que :
\[
f(x_n)\to\ell_1,
\qquad
f(y_n)\to\ell_2,
\]
avec :
\[
\ell_1\neq\ell_2.
\]
\end{methodbox}

\begin{exercicebox}[Non-existence d'une limite]
Montrer que :
\[
f(x)=\sin\left(\frac1x\right)
\]
n'admet pas de limite en \(0\).
\end{exercicebox}

\begin{correctionbox}
On considère les deux suites :
\[
x_n=\frac{1}{\frac{\pi}{2}+2n\pi},
\]
et :
\[
y_n=\frac{1}{-\frac{\pi}{2}+2n\pi}
\]
pour \(n\) assez grand afin que les dénominateurs soient non nuls.

On a :
\[
x_n\to0,
\qquad
y_n\to0.
\]
Mais :
\[
f(x_n)=\sin\left(\frac1{x_n}\right)
=
\sin\left(\frac{\pi}{2}+2n\pi\right)=1,
\]
tandis que :
\[
f(y_n)=\sin\left(\frac1{y_n}\right)
=
\sin\left(-\frac{\pi}{2}+2n\pi\right)=-1.
\]
Les deux limites obtenues sont différentes.  
Donc \(f\) n'admet pas de limite en \(0\).
\end{correctionbox}

%=========================================================
\section{Comparaison et théorème des gendarmes pour les fonctions}

\subsection{Passage à la limite dans les inégalités}

\begin{theorembox}
Soient \(f\) et \(g\) deux fonctions telles que :
\[
f(x)\leq g(x)
\]
au voisinage de \(a\), sauf éventuellement en \(a\).  
Si :
\[
\lim_{x\to a}f(x)=\ell
\quad\text{et}\quad
\lim_{x\to a}g(x)=m,
\]
alors :
\[
\ell\leq m.
\]
\end{theorembox}

\begin{proof}
On peut utiliser la caractérisation séquentielle.  
Soit \((x_n)\) une suite telle que :
\[
x_n\to a,
\qquad x_n\neq a.
\]
Pour \(n\) assez grand, \(x_n\) appartient au voisinage où :
\[
f(x_n)\leq g(x_n).
\]
En passant à la limite dans les suites :
\[
f(x_n)\to \ell,
\qquad
g(x_n)\to m.
\]
On obtient :
\[
\ell\leq m.
\]
\end{proof}

\subsection{Théorème des gendarmes}

\begin{theorembox}[Théorème des gendarmes pour les fonctions]
Soient \(f,g,h\) trois fonctions telles que, au voisinage de \(a\),
\[
f(x)\leq g(x)\leq h(x).
\]
Si :
\[
\lim_{x\to a}f(x)=\ell
\quad\text{et}\quad
\lim_{x\to a}h(x)=\ell,
\]
alors :
\[
\lim_{x\to a}g(x)=\ell.
\]
\end{theorembox}

\begin{proof}
Soit \(\varepsilon>0\).  
Comme \(f(x)\to\ell\), on a, pour \(x\) assez proche de \(a\),
\[
\ell-\varepsilon\leq f(x).
\]
Comme \(h(x)\to\ell\), on a, pour \(x\) assez proche de \(a\),
\[
h(x)\leq \ell+\varepsilon.
\]
Au voisinage de \(a\),
\[
f(x)\leq g(x)\leq h(x).
\]
Donc, pour \(x\) suffisamment proche de \(a\),
\[
\ell-\varepsilon\leq g(x)\leq \ell+\varepsilon.
\]
Ainsi :
\[
|g(x)-\ell|\leq\varepsilon.
\]
Donc :
\[
g(x)\to\ell.
\]
\end{proof}

\begin{exercicebox}[Gendarmes]
Montrer que :
\[
\lim_{x\to0}x\sin\left(\frac1x\right)=0.
\]
\end{exercicebox}

\begin{correctionbox}
Pour tout \(x\neq0\),
\[
-1\leq \sin\left(\frac1x\right)\leq1.
\]
Donc :
\[
-|x|\leq x\sin\left(\frac1x\right)\leq |x|.
\]
Or :
\[
-|x|\to0
\quad\text{et}\quad
|x|\to0
\]
lorsque \(x\to0\).  
Par le théorème des gendarmes :
\[
x\sin\left(\frac1x\right)\to0.
\]
\end{correctionbox}

%=========================================================
\section{Continuité en un point}

\subsection{Définition}

\begin{definitionbox}[Continuité en un point]
Soit \(f:D\to\R\), et soit \(a\in D\).  
La fonction \(f\) est continue en \(a\) si :
\[
\lim_{x\to a}f(x)=f(a).
\]
Autrement dit :
\[
\forall\varepsilon>0,\quad \exists\eta>0,\quad \forall x\in D,\quad
|x-a|\leq\eta \implique |f(x)-f(a)|\leq\varepsilon.
\]
\end{definitionbox}

\begin{retenirbox}
Une fonction est continue en \(a\) lorsque sa limite en \(a\) existe et vaut sa valeur en \(a\).
\[
\lim_{x\to a}f(x)=f(a).
\]
\end{retenirbox}

\subsection{Caractérisation séquentielle de la continuité}

\begin{theorembox}[Caractérisation séquentielle de la continuité]
Soit \(f:D\to\R\) et \(a\in D\).  
La fonction \(f\) est continue en \(a\) si et seulement si, pour toute suite \((x_n)\) d'éléments de \(D\) telle que :
\[
x_n\to a,
\]
on a :
\[
f(x_n)\to f(a).
\]
\end{theorembox}

\begin{proof}
C'est une conséquence directe de la caractérisation séquentielle des limites, avec le fait que \(a\in D\) et que la limite recherchée est \(f(a)\).
\end{proof}

\subsection{Opérations sur les fonctions continues}

\begin{theorembox}[Stabilité par opérations]
Soient \(f\) et \(g\) deux fonctions continues en \(a\).  
Alors :
\[
f+g,\quad fg
\]
sont continues en \(a\).  
Si :
\[
g(a)\neq0,
\]
alors :
\[
\frac{f}{g}
\]
est continue en \(a\).
\end{theorembox}

\begin{proof}
Comme \(f\) et \(g\) sont continues en \(a\),
\[
f(x)\to f(a)
\quad\text{et}\quad
g(x)\to g(a).
\]
Les opérations sur les limites donnent :
\[
f(x)+g(x)\to f(a)+g(a),
\]
\[
f(x)g(x)\to f(a)g(a).
\]
Si \(g(a)\neq0\), alors :
\[
\frac{f(x)}{g(x)}\to\frac{f(a)}{g(a)}.
\]
Cela signifie exactement que les fonctions correspondantes sont continues en \(a\).
\end{proof}

\begin{theorembox}[Continuité d'une composée]
Si \(g\) est continue en \(a\) et si \(f\) est continue en \(g(a)\), alors :
\[
f\circ g
\]
est continue en \(a\).
\end{theorembox}

\begin{proof}
Comme \(g\) est continue en \(a\),
\[
g(x)\to g(a).
\]
Comme \(f\) est continue en \(g(a)\),
\[
f(g(x))\to f(g(a)).
\]
Donc :
\[
f\circ g
\]
est continue en \(a\).
\end{proof}

%=========================================================
\section{Continuité sur un intervalle}

\subsection{Définition}

\begin{definitionbox}[Continuité sur un intervalle]
Soit \(I\) un intervalle de \(\R\).  
Une fonction \(f:I\to\R\) est continue sur \(I\) si elle est continue en tout point de \(I\).
\end{definitionbox}

On note :
\[
f\in\mathcal C^0(I,\R).
\]

\subsection{Exemples fondamentaux}

\begin{theorembox}[Fonctions usuelles continues]
Les fonctions suivantes sont continues sur leur domaine de définition :
\[
x\mapsto x^n,
\qquad
x\mapsto \frac1x,
\qquad
x\mapsto \sqrt{x},
\]
\[
x\mapsto \sin x,
\qquad
x\mapsto \cos x,
\qquad
x\mapsto \tan x,
\]
\[
x\mapsto \exp x,
\qquad
x\mapsto \ln x.
\]
\end{theorembox}

\begin{retenirbox}
Les fonctions obtenues par somme, produit, quotient licite et composition de fonctions usuelles continues sont continues sur leur domaine de définition.
\end{retenirbox}

\begin{exercicebox}[Continuité]
Étudier la continuité de :
\[
f(x)=\frac{\sqrt{x+1}}{x^2-1}.
\]
\end{exercicebox}

\begin{correctionbox}
La fonction est définie lorsque :
\[
x+1\geq0
\]
et :
\[
x^2-1\neq0.
\]
Donc :
\[
x\geq -1,
\]
et :
\[
x\neq -1,\quad x\neq1.
\]
Ainsi :
\[
\Df=]-1,+\infty[\setminus\{1\}.
\]
Sur ce domaine, \(f\) est un quotient de fonctions continues dont le dénominateur ne s'annule pas.  
Donc \(f\) est continue sur :
\[
]-1,1[
\quad\text{et sur}\quad
]1,+\infty[.
\]
\end{correctionbox}

%=========================================================
\section{Prolongement par continuité}

\subsection{Définition}

\begin{definitionbox}[Prolongement par continuité]
Soit \(f:D\to\R\), et soit \(a\notin D\) ou un point où \(f\) n'est pas définie.  
Si :
\[
\lim_{x\to a}f(x)=\ell\in\R,
\]
alors on peut définir une nouvelle fonction \(\widetilde f\) en posant :
\[
\widetilde f(x)=f(x)\quad\text{pour }x\in D,
\]
et :
\[
\widetilde f(a)=\ell.
\]
La fonction \(\widetilde f\) est appelée le prolongement par continuité de \(f\) en \(a\).
\end{definitionbox}

\begin{retenirbox}
Pour prolonger une fonction par continuité en \(a\), il faut que la limite en \(a\) existe et soit finie.
\end{retenirbox}

\begin{exercicebox}[Prolongement par continuité]
Soit :
\[
f(x)=\frac{x^2-1}{x-1}
\]
définie sur \(\R\setminus\{1\}\).  
Montrer que \(f\) est prolongeable par continuité en \(1\).
\end{exercicebox}

\begin{correctionbox}
Pour \(x\neq1\), on factorise :
\[
x^2-1=(x-1)(x+1).
\]
Donc :
\[
f(x)=\frac{(x-1)(x+1)}{x-1}=x+1.
\]
Ainsi :
\[
\lim_{x\to1}f(x)=\lim_{x\to1}(x+1)=2.
\]
On définit :
\[
\widetilde f(1)=2,
\]
et :
\[
\widetilde f(x)=f(x)
\]
pour \(x\neq1\).  
Alors \(\widetilde f\) est continue en \(1\).
\end{correctionbox}

\begin{erreurbox}
Une fonction non définie en \(a\) n'est pas forcément prolongeable par continuité.  
Par exemple :
\[
f(x)=\frac1x
\]
n'est pas prolongeable par continuité en \(0\), car sa limite en \(0\) n'est pas finie.
\end{erreurbox}

%=========================================================
\section{Théorème des valeurs intermédiaires}

\subsection{Énoncé}

\begin{theorembox}[Théorème des valeurs intermédiaires]
Soit \(f:[a,b]\to\R\) une fonction continue sur \([a,b]\).  
Alors, pour tout réel \(y\) compris entre \(f(a)\) et \(f(b)\), il existe \(c\in[a,b]\) tel que :
\[
f(c)=y.
\]
\end{theorembox}

Autrement dit :
\[
[f(a),f(b)]\subset f([a,b])
\]
si \(f(a)\leq f(b)\), et :
\[
[f(b),f(a)]\subset f([a,b])
\]
si \(f(b)\leq f(a)\).

\begin{retenirbox}
Le TVI est un théorème d'existence.  
Il garantit l'existence d'un antécédent, mais pas son unicité.
\end{retenirbox}

\subsection{Démonstration du TVI}

\begin{proof}
On suppose, pour simplifier, que :
\[
f(a)\leq y\leq f(b).
\]
On cherche \(c\in[a,b]\) tel que :
\[
f(c)=y.
\]
Considérons :
\[
A=\{x\in[a,b]: f(x)\leq y\}.
\]
L'ensemble \(A\) est non vide car :
\[
a\in A.
\]
Il est majoré par \(b\).  
Donc il admet une borne supérieure :
\[
c=\sup A.
\]
On veut montrer :
\[
f(c)=y.
\]

Comme \(c\in[a,b]\), la continuité de \(f\) en \(c\) va permettre de conclure.  

Montrons d'abord que :
\[
f(c)\leq y.
\]
Par caractérisation séquentielle de la borne supérieure, il existe une suite \((x_n)\) d'éléments de \(A\) telle que :
\[
x_n\to c.
\]
Pour tout \(n\),
\[
f(x_n)\leq y.
\]
Par continuité de \(f\) en \(c\),
\[
f(x_n)\to f(c).
\]
En passant à la limite dans l'inégalité :
\[
f(c)\leq y.
\]

Montrons maintenant que :
\[
f(c)\geq y.
\]
Si \(c=b\), alors :
\[
f(c)=f(b)\geq y.
\]
Si \(c<b\), pour tout \(n\) assez grand, le réel :
\[
c+\frac1n
\]
appartient à \([a,b]\) et est strictement supérieur à \(c\).  
Comme \(c=\sup A\), aucun réel strictement supérieur à \(c\) suffisamment proche dans \([a,b]\) ne peut appartenir à \(A\).  
Donc :
\[
f\left(c+\frac1n\right)>y
\]
pour \(n\) assez grand.  
En faisant tendre \(n\) vers \(+\infty\), et en utilisant la continuité de \(f\) en \(c\), on obtient :
\[
f(c)\geq y.
\]
Ainsi :
\[
f(c)=y.
\]
\end{proof}

\begin{methodbox}[Utiliser le TVI]
Pour montrer qu'une équation :
\[
f(x)=y
\]
admet au moins une solution sur \([a,b]\), on montre :
\begin{enumerate}
    \item \(f\) est continue sur \([a,b]\) ;
    \item \(y\) est compris entre \(f(a)\) et \(f(b)\).
\end{enumerate}
\end{methodbox}

\begin{exercicebox}[Application du TVI]
Montrer que l'équation :
\[
x^3+x-1=0
\]
admet au moins une solution dans \([0,1]\).
\end{exercicebox}

\begin{correctionbox}
On pose :
\[
f(x)=x^3+x-1.
\]
La fonction \(f\) est polynomiale, donc continue sur \([0,1]\).  
On calcule :
\[
f(0)=-1,
\]
et :
\[
f(1)=1.
\]
Comme :
\[
0\in[-1,1],
\]
le théorème des valeurs intermédiaires assure qu'il existe \(c\in[0,1]\) tel que :
\[
f(c)=0.
\]
Donc l'équation :
\[
x^3+x-1=0
\]
admet au moins une solution dans \([0,1]\).
\end{correctionbox}

%=========================================================
\section{Image d'un intervalle par une fonction continue}

\subsection{Image d'un intervalle}

\begin{theorembox}[Image d'un intervalle par une fonction continue]
Soit \(I\) un intervalle de \(\R\), et soit \(f:I\to\R\) une fonction continue.  
Alors :
\[
f(I)
\]
est un intervalle de \(\R\).
\end{theorembox}

\begin{proof}
Soient \(u,v\in f(I)\) avec :
\[
u\leq v.
\]
Il existe \(a,b\in I\) tels que :
\[
u=f(a),
\qquad
v=f(b).
\]
Soit \(y\in[u,v]\).  
Comme \(I\) est un intervalle, le segment reliant \(a\) et \(b\) est inclus dans \(I\) :
\[
[\min(a,b),\max(a,b)]\subset I.
\]
La fonction \(f\) est continue sur ce segment.  
Par le théorème des valeurs intermédiaires, il existe \(c\) entre \(a\) et \(b\) tel que :
\[
f(c)=y.
\]
Donc :
\[
y\in f(I).
\]
Ainsi, tout réel compris entre deux éléments de \(f(I)\) appartient encore à \(f(I)\).  
Donc \(f(I)\) est un intervalle.
\end{proof}

\begin{retenirbox}
Une fonction continue transforme un intervalle en intervalle.
\end{retenirbox}

\subsection{Image d'un segment}

\begin{theorembox}[Théorème des bornes atteintes]
Soit \(f:[a,b]\to\R\) une fonction continue.  
Alors \(f\) est bornée sur \([a,b]\), et elle atteint ses bornes.  
Autrement dit, il existe \(\alpha,\beta\in[a,b]\) tels que :
\[
f(\alpha)=\min_{x\in[a,b]}f(x),
\]
et :
\[
f(\beta)=\max_{x\in[a,b]}f(x).
\]
\end{theorembox}

\begin{theorembox}[Image continue d'un segment]
Si \(f\) est continue sur \([a,b]\), alors :
\[
f([a,b])
\]
est un segment :
\[
[m,M],
\]
où :
\[
m=\min_{x\in[a,b]}f(x),
\qquad
M=\max_{x\in[a,b]}f(x).
\]
\end{theorembox}

\begin{proof}
Par le théorème des bornes atteintes, \(f\) atteint un minimum \(m\) et un maximum \(M\) sur \([a,b]\).  
Donc :
\[
f([a,b])\subset[m,M].
\]
Comme \(f\) est continue et que \([a,b]\) est un intervalle, \(f([a,b])\) est un intervalle.  
Il contient \(m\) et \(M\).  
Donc il contient tout le segment :
\[
[m,M].
\]
Ainsi :
\[
f([a,b])=[m,M].
\]
\end{proof}

\begin{erreurbox}
L'image d'un intervalle par une fonction quelconque n'est pas forcément un intervalle.  
La continuité est essentielle.
\end{erreurbox}

%=========================================================
\section{Injectivité, surjectivité et bijectivité}

\subsection{Définitions}

\begin{definitionbox}[Injectivité]
Une fonction \(f:E\to F\) est injective si :
\[
\forall x_1,x_2\in E,\quad f(x_1)=f(x_2)\implique x_1=x_2.
\]
\end{definitionbox}

\begin{definitionbox}[Surjectivité]
Une fonction \(f:E\to F\) est surjective si :
\[
\forall y\in F,\quad \exists x\in E,\quad f(x)=y.
\]
\end{definitionbox}

\begin{definitionbox}[Bijectivité]
Une fonction \(f:E\to F\) est bijective si elle est à la fois injective et surjective.  
Autrement dit :
\[
\forall y\in F,\quad \exists! x\in E,\quad f(x)=y.
\]
\end{definitionbox}

\begin{retenirbox}
L'injectivité donne l'unicité d'un antécédent.  
La surjectivité donne l'existence d'un antécédent.  
La bijectivité donne existence et unicité.
\end{retenirbox}

\subsection{Monotonie et injectivité}

\begin{theorembox}
Toute fonction strictement monotone sur un intervalle est injective.
\end{theorembox}

\begin{proof}
Supposons \(f\) strictement croissante.  
Soient \(x_1,x_2\in I\) avec :
\[
x_1\neq x_2.
\]
Si :
\[
x_1<x_2,
\]
alors :
\[
f(x_1)<f(x_2).
\]
Donc :
\[
f(x_1)\neq f(x_2).
\]
Ainsi \(f\) est injective.  
Le cas strictement décroissant est analogue.
\end{proof}

\begin{erreurbox}
Une fonction monotone non strictement monotone n'est pas forcément injective.  
Exemple : une fonction constante est croissante et décroissante, mais pas injective.
\end{erreurbox}

%=========================================================
\section{Théorème de la bijection continue}

\subsection{Énoncé principal}

\begin{theorembox}[Théorème de la bijection continue]
Soit \(I\) un intervalle de \(\R\), et soit \(f:I\to\R\) une fonction continue et strictement monotone.  
Alors \(f\) réalise une bijection de \(I\) sur l'intervalle :
\[
f(I).
\]
De plus, sa fonction réciproque :
\[
f^{-1}:f(I)\to I
\]
est continue et strictement monotone, de même sens de variation que \(f\).
\end{theorembox}

\begin{proof}
Comme \(f\) est strictement monotone, elle est injective.  
Par définition, \(f\) est surjective de \(I\) sur son image \(f(I)\).  
Donc \(f\) est bijective de \(I\) sur \(f(I)\).

Comme \(f\) est continue sur l'intervalle \(I\), l'image \(f(I)\) est un intervalle.  
Il existe donc une fonction réciproque :
\[
f^{-1}:f(I)\to I.
\]

Montrons que \(f^{-1}\) est strictement monotone.  
Si \(f\) est strictement croissante, soient \(y_1<y_2\) dans \(f(I)\).  
On pose :
\[
x_1=f^{-1}(y_1),
\qquad
x_2=f^{-1}(y_2).
\]
Si \(x_1\geq x_2\), alors, comme \(f\) est strictement croissante :
\[
f(x_1)\geq f(x_2),
\]
donc :
\[
y_1\geq y_2,
\]
contradiction.  
Donc :
\[
x_1<x_2.
\]
Ainsi \(f^{-1}\) est strictement croissante.

Le cas décroissant est analogue : la réciproque est strictement décroissante.

La continuité de \(f^{-1}\) est un résultat fondamental qui découle du fait qu'une fonction continue strictement monotone transforme les intervalles en intervalles sans saut.  
On peut le démontrer par caractérisation séquentielle : si \(y_n\to y\) dans \(f(I)\), alors les \(x_n=f^{-1}(y_n)\) ne peuvent pas avoir de limite autre que \(x=f^{-1}(y)\), sinon la continuité et la stricte monotonie de \(f\) donneraient une contradiction.  
Ainsi :
\[
f^{-1}(y_n)\to f^{-1}(y),
\]
ce qui prouve la continuité de \(f^{-1}\).
\end{proof}

\begin{methodbox}[Appliquer le théorème de la bijection continue]
Pour montrer qu'une équation :
\[
f(x)=y
\]
admet une unique solution sur un intervalle \(I\), on montre :
\begin{enumerate}
    \item \(f\) est continue sur \(I\) ;
    \item \(f\) est strictement monotone sur \(I\) ;
    \item \(y\in f(I)\).
\end{enumerate}
Alors il existe un unique \(x\in I\) tel que :
\[
f(x)=y.
\]
\end{methodbox}

\begin{exercicebox}[Bijection continue]
Montrer que la fonction :
\[
f:x\mapsto x^3+x
\]
réalise une bijection de \(\R\) sur \(\R\).
\end{exercicebox}

\begin{correctionbox}
La fonction \(f\) est polynomiale, donc continue sur \(\R\).

Soient \(x<y\).  
Alors :
\[
f(y)-f(x)=y^3-x^3+y-x.
\]
On factorise :
\[
y^3-x^3=(y-x)(y^2+xy+x^2).
\]
Donc :
\[
f(y)-f(x)=(y-x)(y^2+xy+x^2+1).
\]
Or :
\[
y^2+xy+x^2
=
\frac12(x^2+y^2)+\frac12(x+y)^2
\geq0.
\]
Donc :
\[
y^2+xy+x^2+1>0.
\]
Comme \(y-x>0\), on obtient :
\[
f(y)-f(x)>0.
\]
Ainsi \(f\) est strictement croissante.

Enfin :
\[
\lim_{x\to-\infty}(x^3+x)=-\infty,
\]
et :
\[
\lim_{x\to+\infty}(x^3+x)=+\infty.
\]
Donc :
\[
f(\R)=\R.
\]
Par le théorème de la bijection continue, \(f\) réalise une bijection de \(\R\) sur \(\R\).
\end{correctionbox}

%=========================================================
\section{Fonction réciproque}

\subsection{Définition}

\begin{definitionbox}[Fonction réciproque]
Soit \(f:E\to F\) une bijection.  
La fonction réciproque de \(f\), notée :
\[
f^{-1}:F\to E,
\]
est définie par :
\[
f^{-1}(y)=x
\]
où \(x\) est l'unique élément de \(E\) tel que :
\[
f(x)=y.
\]
\end{definitionbox}

On a :
\[
f^{-1}\circ f=\Id_E,
\]
et :
\[
f\circ f^{-1}=\Id_F.
\]

\begin{retenirbox}
L'expression :
\[
f^{-1}(y)
\]
désigne l'antécédent de \(y\) par \(f\), lorsque \(f\) est bijective.
\end{retenirbox}

\subsection{Graphe de la réciproque}

Si \(f:I\to J\) est une bijection, alors les graphes de \(f\) et de \(f^{-1}\) sont symétriques par rapport à la droite :
\[
y=x.
\]

Cette propriété est géométrique, mais elle se démontre algébriquement :
\[
y=f(x)
\ssi
x=f^{-1}(y).
\]

\begin{exercicebox}[Réciproque]
Soit :
\[
f:[0,+\infty[\to[0,+\infty[,
\qquad
f(x)=x^2.
\]
Montrer que \(f\) est bijective et déterminer sa fonction réciproque.
\end{exercicebox}

\begin{correctionbox}
La fonction \(x\mapsto x^2\) est continue et strictement croissante sur \([0,+\infty[\).  
De plus :
\[
f(0)=0,
\]
et :
\[
\lim_{x\to+\infty}x^2=+\infty.
\]
Donc :
\[
f([0,+\infty[)=[0,+\infty[.
\]
Par le théorème de la bijection continue, \(f\) est une bijection de \([0,+\infty[\) sur \([0,+\infty[\).

Pour \(y\geq0\), l'équation :
\[
x^2=y
\]
avec \(x\geq0\) admet pour unique solution :
\[
x=\sqrt y.
\]
Donc :
\[
f^{-1}(y)=\sqrt y.
\]
\end{correctionbox}

%=========================================================
\newpage
\section{Exercices progressifs}

\subsection*{Thème 1 -- Domaine, image, parité}

\begin{exercicebox}[1]
Déterminer le domaine de définition de :
\[
f(x)=\frac{\sqrt{2x-1}}{x^2-4}.
\]
\end{exercicebox}

\begin{exercicebox}[2]
Soit :
\[
f(x)=x^2-2x.
\]
Déterminer :
\[
f([0,3]).
\]
\end{exercicebox}

\begin{exercicebox}[3]
Étudier la parité de :
\[
f(x)=\frac{x^3}{1+x^2}.
\]
\end{exercicebox}

\subsection*{Thème 2 -- Limites}

\begin{exercicebox}[4]
Calculer :
\[
\lim_{x\to2}\frac{x^2-4}{x-2}.
\]
\end{exercicebox}

\begin{exercicebox}[5]
Calculer les limites à gauche et à droite en \(0\) de :
\[
f(x)=\frac{|x|}{x}.
\]
\end{exercicebox}

\begin{exercicebox}[6]
Montrer que :
\[
\lim_{x\to0}x^2\cos\left(\frac1x\right)=0.
\]
\end{exercicebox}

\subsection*{Thème 3 -- Continuité}

\begin{exercicebox}[7]
Étudier la continuité de :
\[
f(x)=
\begin{cases}
\frac{x^2-1}{x-1} & \text{si }x\neq1,\\
3 & \text{si }x=1.
\end{cases}
\]
\end{exercicebox}

\begin{exercicebox}[8]
Déterminer \(a\in\R\) pour que la fonction :
\[
f(x)=
\begin{cases}
\frac{\sin x}{x} & \text{si }x\neq0,\\
a & \text{si }x=0
\end{cases}
\]
soit continue en \(0\).
\end{exercicebox}

\subsection*{Thème 4 -- TVI et bijection}

\begin{exercicebox}[9]
Montrer que l'équation :
\[
\cos x=x
\]
admet au moins une solution dans \([0,1]\).
\end{exercicebox}

\begin{exercicebox}[10]
Montrer que l'équation :
\[
x^5+x-1=0
\]
admet une unique solution réelle.
\end{exercicebox}

\subsection*{Thème 5 -- Fonction réciproque}

\begin{exercicebox}[11]
Montrer que :
\[
f:x\mapsto e^x
\]
réalise une bijection de \(\R\) sur \(]0,+\infty[\).  
Identifier sa fonction réciproque.
\end{exercicebox}

\begin{exercicebox}[12]
Soit :
\[
f(x)=x+\ln x
\]
définie sur \(]0,+\infty[\).  
Montrer que \(f\) réalise une bijection de \(]0,+\infty[\) sur \(\R\).
\end{exercicebox}

%=========================================================
\newpage
\section{Corrigés détaillés des exercices progressifs}

\begin{correctionbox}[1]
On considère :
\[
f(x)=\frac{\sqrt{2x-1}}{x^2-4}.
\]
On doit avoir :
\[
2x-1\geq0,
\]
donc :
\[
x\geq\frac12.
\]
On doit aussi avoir :
\[
x^2-4\neq0,
\]
donc :
\[
x\neq -2,\qquad x\neq2.
\]
Comme \(x\geq\frac12\), la condition \(x\neq -2\) est automatiquement vérifiée.  
Donc :
\[
\Df=\left[\frac12,+\infty\right[\setminus\{2\}.
\]
\end{correctionbox}

\begin{correctionbox}[2]
On a :
\[
f(x)=x^2-2x=(x-1)^2-1.
\]
Sur \([0,3]\), le minimum est atteint en :
\[
x=1,
\]
et vaut :
\[
f(1)=-1.
\]
Aux bornes :
\[
f(0)=0,
\qquad
f(3)=3.
\]
Le maximum sur \([0,3]\) vaut donc \(3\).  
Comme \(f\) est continue sur \([0,3]\), son image est un intervalle.  
Donc :
\[
f([0,3])=[-1,3].
\]
\end{correctionbox}

\begin{correctionbox}[3]
Le domaine de définition est \(\R\), symétrique par rapport à \(0\).  
On calcule :
\[
f(-x)=\frac{(-x)^3}{1+(-x)^2}
=
-\frac{x^3}{1+x^2}
=
-f(x).
\]
Donc \(f\) est impaire.
\end{correctionbox}

\begin{correctionbox}[4]
Pour \(x\neq2\),
\[
\frac{x^2-4}{x-2}
=
\frac{(x-2)(x+2)}{x-2}
=
x+2.
\]
Donc :
\[
\lim_{x\to2}\frac{x^2-4}{x-2}
=
\lim_{x\to2}(x+2)
=
4.
\]
\end{correctionbox}

\begin{correctionbox}[5]
Si \(x>0\), alors :
\[
|x|=x.
\]
Donc :
\[
\frac{|x|}{x}=1.
\]
Ainsi :
\[
\lim_{x\to0^+}\frac{|x|}{x}=1.
\]
Si \(x<0\), alors :
\[
|x|=-x.
\]
Donc :
\[
\frac{|x|}{x}=-1.
\]
Ainsi :
\[
\lim_{x\to0^-}\frac{|x|}{x}=-1.
\]
Les deux limites latérales sont différentes.  
Donc la limite en \(0\) n'existe pas.
\end{correctionbox}

\begin{correctionbox}[6]
Pour tout \(x\neq0\),
\[
-1\leq \cos\left(\frac1x\right)\leq1.
\]
Donc :
\[
-x^2\leq x^2\cos\left(\frac1x\right)\leq x^2.
\]
Or :
\[
-x^2\to0,
\qquad
x^2\to0
\]
lorsque \(x\to0\).  
Par le théorème des gendarmes :
\[
x^2\cos\left(\frac1x\right)\to0.
\]
\end{correctionbox}

\begin{correctionbox}[7]
Pour \(x\neq1\),
\[
\frac{x^2-1}{x-1}
=
x+1.
\]
Donc :
\[
\lim_{x\to1}\frac{x^2-1}{x-1}=2.
\]
Or :
\[
f(1)=3.
\]
Comme :
\[
\lim_{x\to1}f(x)\neq f(1),
\]
la fonction n'est pas continue en \(1\).
\end{correctionbox}

\begin{correctionbox}[8]
On sait que :
\[
\lim_{x\to0}\frac{\sin x}{x}=1.
\]
Pour que \(f\) soit continue en \(0\), il faut et il suffit que :
\[
a=1.
\]
Donc la valeur cherchée est :
\[
a=1.
\]
\end{correctionbox}

\begin{correctionbox}[9]
On pose :
\[
f(x)=\cos x-x.
\]
La fonction \(f\) est continue sur \([0,1]\).  
On a :
\[
f(0)=1>0.
\]
Et :
\[
f(1)=\cos1-1<0.
\]
Comme \(f\) change de signe entre \(0\) et \(1\), le TVI assure qu'il existe \(c\in[0,1]\) tel que :
\[
f(c)=0.
\]
Donc :
\[
\cos c=c.
\]
\end{correctionbox}

\begin{correctionbox}[10]
On pose :
\[
f(x)=x^5+x-1.
\]
La fonction \(f\) est polynomiale, donc continue sur \(\R\).  
De plus :
\[
\lim_{x\to-\infty}f(x)=-\infty,
\qquad
\lim_{x\to+\infty}f(x)=+\infty.
\]
Donc, par le TVI, l'équation :
\[
f(x)=0
\]
admet au moins une solution réelle.

Montrons l'unicité.  
Soient \(x<y\).  
Alors :
\[
f(y)-f(x)=y^5-x^5+y-x.
\]
On a :
\[
y-x>0.
\]
De plus, la fonction \(x\mapsto x^5+x\) est strictement croissante, car si \(x<y\), alors \(y^5>x^5\) et \(y>x\).  
Donc :
\[
f(y)>f(x).
\]
Ainsi \(f\) est strictement croissante, donc injective.  
L'équation \(f(x)=0\) admet au plus une solution.  
Elle admet donc exactement une solution réelle.
\end{correctionbox}

\begin{correctionbox}[11]
La fonction exponentielle est continue et strictement croissante sur \(\R\).  
On a :
\[
\lim_{x\to-\infty}e^x=0,
\qquad
\lim_{x\to+\infty}e^x=+\infty.
\]
Donc :
\[
\exp(\R)=]0,+\infty[.
\]
Par le théorème de la bijection continue, \(\exp\) réalise une bijection :
\[
\R\to]0,+\infty[.
\]
Sa fonction réciproque est le logarithme népérien :
\[
\ln: ]0,+\infty[\to\R.
\]
\end{correctionbox}

\begin{correctionbox}[12]
On considère :
\[
f(x)=x+\ln x
\]
sur \(]0,+\infty[\).  
La fonction est continue sur \(]0,+\infty[\).  
Elle est strictement croissante car \(x\mapsto x\) et \(x\mapsto\ln x\) sont strictement croissantes sur cet intervalle.

De plus :
\[
\lim_{x\to0^+}(x+\ln x)=-\infty,
\]
et :
\[
\lim_{x\to+\infty}(x+\ln x)=+\infty.
\]
Donc :
\[
f(]0,+\infty[)=\R.
\]
Par le théorème de la bijection continue, \(f\) réalise une bijection de \(]0,+\infty[\) sur \(\R\).
\end{correctionbox}

%=========================================================
\newpage
\section{Grand devoir bilan final}

\begin{center}
\Large\textbf{Devoir bilan -- Fonctions réelles : limites, continuité et bijections}
\end{center}

\subsection*{Exercice 1 -- Questions de cours}

\begin{enumerate}
    \item Définir la limite finie d'une fonction \(f\) en un réel \(a\).
    \item Définir la continuité de \(f\) en \(a\).
    \item Énoncer le théorème des valeurs intermédiaires.
    \item Énoncer le théorème de la bijection continue.
    \item Expliquer la différence entre image réciproque d'un ensemble et fonction réciproque.
\end{enumerate}

\subsection*{Exercice 2 -- Limites}

Calculer, lorsqu'elles existent :
\[
\lim_{x\to1}\frac{x^3-1}{x-1},
\]
\[
\lim_{x\to0}x\sin\left(\frac1{x^2}\right),
\]
\[
\lim_{x\to0^+}\ln x.
\]

\subsection*{Exercice 3 -- Continuité et prolongement}

Soit :
\[
f(x)=\frac{\sqrt{x+1}-1}{x}
\]
définie pour \(x>-1\) et \(x\neq0\).

\begin{enumerate}
    \item Calculer :
    \[
    \lim_{x\to0}f(x).
    \]
    \item En déduire que \(f\) est prolongeable par continuité en \(0\).
\end{enumerate}

\subsection*{Exercice 4 -- TVI et unicité}

Montrer que l'équation :
\[
x^3+2x-2=0
\]
admet une unique solution réelle.

\subsection*{Exercice 5 -- Bijection réciproque}

Soit :
\[
f:[0,+\infty[\to[1,+\infty[,
\qquad
f(x)=x^2+1.
\]

\begin{enumerate}
    \item Montrer que \(f\) est bijective.
    \item Déterminer \(f^{-1}\).
    \item Étudier la monotonie de \(f^{-1}\).
\end{enumerate}

\newpage

%=========================================================
\section{Correction détaillée du devoir bilan}

\subsection*{Correction de l'exercice 1}

\textbf{1.}  
On dit que :
\[
\lim_{x\to a}f(x)=\ell
\]
si :
\[
\forall\varepsilon>0,\quad \exists\eta>0,\quad
0<|x-a|\leq\eta\implique |f(x)-\ell|\leq\varepsilon.
\]

\textbf{2.}  
La fonction \(f\) est continue en \(a\in D_f\) si :
\[
\lim_{x\to a}f(x)=f(a).
\]

\textbf{3.}  
Si \(f\) est continue sur \([a,b]\), alors tout réel compris entre \(f(a)\) et \(f(b)\) possède au moins un antécédent dans \([a,b]\).

\textbf{4.}  
Si \(f\) est continue et strictement monotone sur un intervalle \(I\), alors \(f\) réalise une bijection de \(I\) sur \(f(I)\), et sa réciproque est continue et strictement monotone.

\textbf{5.}  
L'image réciproque :
\[
f^{-1}(B)
\]
d'un ensemble \(B\) existe pour toute fonction \(f\).  
Elle désigne :
\[
\{x:f(x)\in B\}.
\]
La fonction réciproque \(f^{-1}\) n'existe que lorsque \(f\) est bijective.

\subsection*{Correction de l'exercice 2}

Pour \(x\neq1\),
\[
\frac{x^3-1}{x-1}
=
\frac{(x-1)(x^2+x+1)}{x-1}
=
x^2+x+1.
\]
Donc :
\[
\lim_{x\to1}\frac{x^3-1}{x-1}
=
1+1+1=3.
\]

Ensuite :
\[
-1\leq \sin\left(\frac1{x^2}\right)\leq1.
\]
Donc :
\[
-|x|\leq x\sin\left(\frac1{x^2}\right)\leq |x|.
\]
Par encadrement :
\[
\lim_{x\to0}x\sin\left(\frac1{x^2}\right)=0.
\]

Enfin :
\[
\lim_{x\to0^+}\ln x=-\infty.
\]

\subsection*{Correction de l'exercice 3}

On considère :
\[
f(x)=\frac{\sqrt{x+1}-1}{x}.
\]
On multiplie par la quantité conjuguée :
\[
f(x)
=
\frac{\sqrt{x+1}-1}{x}\cdot
\frac{\sqrt{x+1}+1}{\sqrt{x+1}+1}.
\]
Donc :
\[
f(x)
=
\frac{x+1-1}{x(\sqrt{x+1}+1)}
=
\frac{x}{x(\sqrt{x+1}+1)}.
\]
Pour \(x\neq0\),
\[
f(x)=\frac1{\sqrt{x+1}+1}.
\]
Ainsi :
\[
\lim_{x\to0}f(x)
=
\frac1{\sqrt1+1}
=
\frac12.
\]
Donc \(f\) est prolongeable par continuité en \(0\) en posant :
\[
\widetilde f(0)=\frac12.
\]

\subsection*{Correction de l'exercice 4}

On pose :
\[
f(x)=x^3+2x-2.
\]
La fonction \(f\) est continue sur \(\R\).  
On calcule :
\[
f(0)=-2<0,
\]
et :
\[
f(1)=1>0.
\]
Par le TVI, il existe au moins une solution dans \([0,1]\).

Montrons l'unicité.  
Soient \(x<y\).  
Alors :
\[
f(y)-f(x)=y^3-x^3+2(y-x).
\]
On a :
\[
y^3-x^3=(y-x)(y^2+xy+x^2).
\]
Donc :
\[
f(y)-f(x)
=
(y-x)(y^2+xy+x^2+2).
\]
Or :
\[
y^2+xy+x^2
=
\frac12(x^2+y^2)+\frac12(x+y)^2
\geq0.
\]
Donc :
\[
y^2+xy+x^2+2>0.
\]
Comme \(y-x>0\), on obtient :
\[
f(y)-f(x)>0.
\]
Donc \(f\) est strictement croissante.  
Elle est donc injective.  
L'équation :
\[
f(x)=0
\]
admet au plus une solution.  
Elle en admet au moins une.  
Donc elle admet une unique solution réelle.
\]

\subsection*{Correction de l'exercice 5}

La fonction :
\[
f(x)=x^2+1
\]
est continue sur \([0,+\infty[\).  
Elle est strictement croissante sur \([0,+\infty[\).  
De plus :
\[
f(0)=1,
\]
et :
\[
\lim_{x\to+\infty}f(x)=+\infty.
\]
Donc :
\[
f([0,+\infty[)=[1,+\infty[.
\]
Par le théorème de la bijection continue, \(f\) est bijective de \([0,+\infty[\) sur \([1,+\infty[\).

Pour \(y\in[1,+\infty[\), on résout :
\[
y=x^2+1.
\]
Donc :
\[
x^2=y-1.
\]
Comme \(x\geq0\), on obtient :
\[
x=\sqrt{y-1}.
\]
Ainsi :
\[
f^{-1}(y)=\sqrt{y-1}.
\]
La fonction \(f^{-1}\) est strictement croissante sur \([1,+\infty[\), comme \(f\).

%=========================================================
\newpage
\section{Fiche de synthèse finale}

\subsection*{1. Fonction réelle}

Une fonction réelle d'une variable réelle est une application :
\[
f:D\subset\R\to\R.
\]

Son domaine de définition est l'ensemble des réels pour lesquels \(f(x)\) a un sens.

\subsection*{2. Image directe et image réciproque}

\[
f(A)=\{f(x):x\in A\}.
\]

\[
f^{-1}(B)=\{x\in D:f(x)\in B\}.
\]

L'image réciproque d'un ensemble existe toujours.  
La fonction réciproque n'existe que pour une bijection.

\subsection*{3. Parité et périodicité}

Fonction paire :
\[
f(-x)=f(x).
\]

Fonction impaire :
\[
f(-x)=-f(x).
\]

Fonction périodique :
\[
f(x+T)=f(x).
\]

\subsection*{4. Limite finie en un point}

\[
\lim_{x\to a}f(x)=\ell
\]
signifie :
\[
\forall\varepsilon>0,\quad \exists\eta>0,\quad
0<|x-a|\leq\eta\implique |f(x)-\ell|\leq\varepsilon.
\]

\subsection*{5. Limites latérales}

\[
\lim_{x\to a}f(x)=\ell
\]
si et seulement si :
\[
\lim_{x\to a^-}f(x)=\ell
\quad\text{et}\quad
\lim_{x\to a^+}f(x)=\ell.
\]

\subsection*{6. Caractérisation séquentielle}

\[
\lim_{x\to a}f(x)=\ell
\]
si et seulement si, pour toute suite \(x_n\to a\), avec \(x_n\neq a\), on a :
\[
f(x_n)\to\ell.
\]

\subsection*{7. Continuité}

\(f\) est continue en \(a\) si :
\[
\lim_{x\to a}f(x)=f(a).
\]

\(f\) est continue sur \(I\) si elle est continue en tout point de \(I\).

\subsection*{8. Prolongement par continuité}

Si :
\[
\lim_{x\to a}f(x)=\ell\in\R,
\]
alors on peut prolonger \(f\) par continuité en posant :
\[
\widetilde f(a)=\ell.
\]

\subsection*{9. Théorème des valeurs intermédiaires}

Si \(f\) est continue sur \([a,b]\), alors tout réel compris entre \(f(a)\) et \(f(b)\) possède un antécédent dans \([a,b]\).

\subsection*{10. Image d'un intervalle}

L'image d'un intervalle par une fonction continue est un intervalle.

L'image d'un segment par une fonction continue est un segment :
\[
f([a,b])=[m,M].
\]

\subsection*{11. Théorème de la bijection continue}

Si \(f\) est continue et strictement monotone sur un intervalle \(I\), alors :
\[
f:I\to f(I)
\]
est bijective.

Sa réciproque :
\[
f^{-1}:f(I)\to I
\]
est continue et strictement monotone.

\subsection*{12. Méthodes essentielles}

Pour montrer l'existence d'une solution :
\[
f(x)=0,
\]
on utilise souvent le TVI.

Pour montrer l'existence et l'unicité :
\[
f(x)=0,
\]
on combine :
\[
\text{continuité}
+
\text{changement de signe}
+
\text{stricte monotonie}.
\]

\begin{center}
\[
\boxed{
\text{Continuité donne l'existence ; stricte monotonie donne l'unicité.}
}
\]
\end{center}

%=========================================================
\section*{Conclusion pédagogique}
\addcontentsline{toc}{section}{Conclusion pédagogique}

Ce chapitre établit les bases de l'analyse des fonctions réelles.  
La notion de limite permet de décrire le comportement local d'une fonction.  
La continuité traduit l'absence de rupture.  
Le théorème des valeurs intermédiaires exprime une propriété profonde des fonctions continues sur un intervalle : elles ne sautent aucune valeur intermédiaire.

Le théorème de la bijection continue est l'un des outils les plus utilisés en prépa.  
Il permet de montrer proprement l'existence et l'unicité de solutions, ainsi que de définir rigoureusement des fonctions réciproques.

\begin{center}
\[
\boxed{
\text{Fonctions réelles : comprendre les variations, les limites, et les valeurs prises.}
}
\]
\end{center}

\end{document}